% Algorithmic Non-commutative Class Field Theory, III
% The Iwahori-Hecke tower and the type III dissolution of the trace obstruction
% Build: pdflatex anccft3 (x3)
\documentclass[11pt]{amsart}
\usepackage{amsmath,amssymb,amsthm,mathtools}
\usepackage[margin=1.0in]{geometry}
\usepackage{booktabs}
\raggedbottom
\usepackage[hidelinks]{hyperref}

\numberwithin{equation}{section}
\newtheorem{theorem}{Theorem}[section]
\newtheorem{proposition}[theorem]{Proposition}
\newtheorem{lemma}[theorem]{Lemma}
\newtheorem{corollary}[theorem]{Corollary}
\newtheorem{hypothesis}[theorem]{Hypothesis}
\theoremstyle{definition}
\newtheorem{definition}[theorem]{Definition}
\newtheorem{conjecture}[theorem]{Conjecture}
\theoremstyle{remark}
\newtheorem{remark}[theorem]{Remark}

\newcommand{\Z}{\mathbb{Z}}
\newcommand{\Q}{\mathbb{Q}}
\newcommand{\R}{\mathbb{R}}
\newcommand{\C}{\mathbb{C}}
\newcommand{\N}{\mathbb{N}}
\newcommand{\cE}{\mathcal{E}}
\newcommand{\cG}{\mathcal{G}}
\newcommand{\cH}{\mathcal{H}}
\newcommand{\cI}{\mathcal{I}}
\newcommand{\cR}{\mathcal{R}}
\newcommand{\Tr}{\operatorname{Tr}}
\newcommand{\Jac}{\operatorname{Jac}}
\newcommand{\coker}{\operatorname{coker}}
\newcommand{\GL}{\mathrm{GL}}
\newcommand{\PGL}{\mathrm{PGL}}
\newcommand{\ord}{\operatorname{ord}}
\newcommand{\supp}{\operatorname{supp}}

\title[Algorithmic non-commutative class field theory, III]
{Algorithmic non-commutative class field theory, III:\\
the Iwahori--Hecke tower and the type III dissolution\\ of the trace obstruction}

\author{Ruqing Chen}
\address{GUT Geoservice Inc., Rivi\`ere-Beaudette, Qu\'ebec, Canada}
\email{ruqing@hotmail.com}
\date{8 September 2026}
\subjclass[2020]{Primary 46L36, 46L55; Secondary 37A20, 46L80, 11F70, 05C50}
\keywords{Krieger factor, asymptotic ratio set, type III factor, tail equivalence,
Bruhat--Tits tree, Iwahori--Hecke algebra, sandpile group, Iwasawa invariants}

\thanks{Paper III of the series. Papers I--II are cited internally as
[I] and [II]; their numbered results are quoted with those prefixes. External
citations are subject to the series' standing caveat --- bibliographic data from
recollection and out-of-band checks --- except where an entry is explicitly marked
as verified against its primary source.}

\begin{document}

\begin{abstract}
Papers I--II established that the arithmetic of the $\GL_2$ Hecke--Laplacian
$\Delta_p=(q+1)I-T_p$ --- the critical group $\Jac(Y)$ of order
$\frac1n\prod_f\#E_f(\mathbb F_p)$ --- is carried by a virtual $KK$-class whose
companion matrix $C$ satisfies $\Tr(C^2)=-n(q-1)<0$, so that no non-negative
$C^*$-correspondence of any size realizes it (the unconditional trace obstruction,
[II, Thm.~1.6]). This paper takes the projective limit over the tower and determines
what survives.

Section~\ref{sec:type} proves that the geometric tail equivalence relation on the
Cantor limit $\Omega=\varprojlim_kV_{Y_k}$, with its harmonic measure $\nu$, has
Radon--Nikodym cocycle valued in $q^{\Z}$ and asymptotic ratio set
$r_\infty=q^{\Z}\cup\{0\}$, computed in full by a density argument; by
Krieger's theorem the associated von Neumann algebra is the hyperfinite factor of type
$\mathrm{III}_{1/q}$, the Powers factor $R_{1/q}$, and for Satake-twisted cocycle
classes with range $\mu_0^{\Z}$ the same computation gives type
$\mathrm{III}_{\mu_0\wedge\mu_0^{-1}}$. A type III factor carries no faithful normal
semifinite trace, so the inequality $\Tr(\widetilde C^2)\ge0$ that powered the finite
No-Go theorems has no formulation on the limit: the obstruction dissolves as a
\emph{sentence}, not as a defeated inequality. We prove equally that the dissolution
buys less than it appears to: any inductive-limit realization by finite-level
correspondences still dies level by level, and the weak closure that dissolves the
trace also annihilates the $K$-theory the programme wants to realize --- the correct
target is a traceless Kirchberg $C^*$-limit, not a von Neumann algebra, and we say so.

Section~\ref{sec:euler} formalizes the finding of [II, Rem.~5.2]: the congruence
(Iwahori path) tower obeys the exact pseudo-Iwasawa law $\nu_2=6\cdot2^k-k-4$ with
$\lambda=-1$, impossible for a $\Z_p[[T]]$-characteristic ideal. The one-step
recursion $|\mathrm{sandpile}(Y_{k+1})|=q^{\,n_k(q-1)-1}|\mathrm{sandpile}(Y_k)|$ ---
whose exponent corrects the $q=2$ special form of earlier drafts --- is proved here
outright, self-containedly, from the $AB$/$BA$ spectral identity for line digraphs
together with the equal-cofactor eigenvalue-product formula for Eulerian Laplacians,
making Section~\ref{sec:euler} unconditional; the full law follows in two lines, and
$\lambda=-1$ is identified as $-\chi$ for the one-dimensional Eisenstein kernel removed
at each level of the degeneracy correspondences: an Euler-characteristic correction
over the Iwahori--Hecke base, external to any commutative Iwasawa module. A closing
outlook section states, without derivation, the $p$-adic horizon.
\end{abstract}

\maketitle

\section*{Status of the results}

Theorem~\ref{thm:type} and its supporting lemmas are proved here in full up to two
classical inputs, quoted and flagged: Krieger's ratio-set criterion for the type of an
ergodic amenable relation, and the uniqueness of the hyperfinite type
$\mathrm{III}_\lambda$ factor for $0<\lambda<1$. Corollary~\ref{cor:dissolve} and
Proposition~\ref{prop:inductive} are complete. Section~\ref{sec:euler} is now
unconditional: the line-digraph recursion, previously a five-level computational
hypothesis, is proved in full as Theorem~\ref{thm:recursion}, self-containedly and
with the exponent corrected from $n_k-1$ (a $q=2$ artifact) to $n_k(q-1)-1$; the
computational confirmations ($q=2$, five levels; $q=3$, one instance) now confirm a
theorem rather than support a hypothesis. Section~\ref{sec:outlook} derives nothing
beyond one exactly verified bookkeeping identity, clearly labelled.

\section{The type of the limit, and what dissolves}\label{sec:type}

\subsection{The Cantor limit and its harmonic measure}

Throughout, $q\ge2$. Let $X$ be the $(q+1)$-regular tree with base vertex $x_0$ and let
\[
\Omega\;=\;\{1,\dots,q+1\}\times\{1,\dots,q\}^{\N}
\]
be its geodesic boundary in rooted coordinates: the first symbol chooses among the
$q+1$ edges at $x_0$, every later symbol among the $q$ non-backtracking continuations.
Under the identifications of [I, \S1], $\Omega\cong\partial X\cong\varprojlim_kV_{Y_k}$
for the Iwahori path tower, a Cantor set. The \emph{harmonic measure} $\nu$ is the
product measure with uniform factors: mass $\tfrac1{q+1}$ on the first coordinate,
$\tfrac1q$ on each subsequent one, so that every cylinder $[w]$ of length $n\ge1$ has
$\nu([w])=\tfrac1{q+1}q^{-(n-1)}$.

\begin{definition}[Tail relation and degree cocycle]
$\cR$ is the tail equivalence relation on $\Omega$:
$\omega\sim\omega'$ iff $\sigma^m\omega=\sigma^n\omega'$ for some $m,n\ge0$, where
$\sigma$ deletes the first symbol. The \emph{degree cocycle} of the pair
$(\omega,\omega')$ realized at $(m,n)$ is $c=m-n\in\Z$; it is well defined on $\cR$.
$\cR_0\subset\cR$ is the degree-zero subrelation.
\end{definition}

\begin{lemma}[Radon--Nikodym cocycle]\label{lem:rn}
$\nu$ is quasi-invariant under $\cR$, and every partial isomorphism
$\gamma\in[[\cR]]$ of degree $c(\gamma)$ has
\[
\frac{d\nu\circ\gamma}{d\nu}=q^{\,c(\gamma)}\qquad\nu\text{-a.e.\ on its domain.}
\]
In particular the cocycle takes values in the discrete group $q^{\Z}\subset\R_{>0}$,
and $\cR_0$ preserves $\nu$.
\end{lemma}

\begin{proof}
It suffices to check the generating partial isomorphisms
$T_{u,v}\colon[u]\to[v]$, $(u\xi)\mapsto(v\xi)$, for finite words $u,v$ of lengths
$m,n\ge1$: these have degree $m-n$ and scale every cylinder inside $[u]$ by
$\nu([v])/\nu([u])=q^{\,m-n}$. Every element of $[[\cR]]$ is a countable disjoint union
of restrictions of such maps.
\end{proof}

\begin{lemma}[Ergodicity]\label{lem:ergodic}
$\cR_0$, hence $\cR$, is ergodic on $(\Omega,\nu)$.
\end{lemma}

\begin{proof}
A Borel set invariant under $\cR_0$ is invariant under every change of finitely many
initial coordinates (compose $T_{u,v}$ with $|u|=|v|$), hence is a tail event for the
product measure $\nu$; Kolmogorov's zero--one law gives measure $0$ or $1$.
\end{proof}

\subsection{The ratio set}

Recall Krieger's asymptotic ratio set $r_\infty(\cR,\nu)\subseteq[0,\infty)$:
$\lambda\in r_\infty$ iff for every $\varepsilon>0$ and every Borel $A$ with
$\nu(A)>0$ there is $\gamma\in[[\cR]]$ with domain $B\subseteq A$ of positive measure,
$\gamma(B)\subseteq A$, and $|d\nu\circ\gamma/d\nu-\lambda|<\varepsilon$ on $B$.

\begin{theorem}[Ratio set of the geometric tail relation]\label{thm:ratio}
$r_\infty(\cR,\nu)=q^{\Z}\cup\{0\}$.
\end{theorem}

\begin{proof}
\emph{Upper bound.} By Lemma~\ref{lem:rn} the derivative of every partial isomorphism
is $\nu$-a.e.\ $q^{\Z}$-valued, and $q^{\Z}$ is closed and discrete in $\R_{>0}$; hence
$r_\infty\cap\R_{>0}\subseteq q^{\Z}$.

\emph{Lower bound: $q\in r_\infty$.} Let $A$ have positive measure and let
$0<\delta<\tfrac13$. By the martingale (Lebesgue density) theorem for the cylinder
filtration, there is a point of $A$ at which the cylinder densities converge to $1$;
choose nested cylinders $[u]\supset[us]$ (with $|u|\ge1$, $s$ the next symbol of that
point) satisfying
\[
\nu(A\cap[u])>(1-\delta)\,\nu([u]),\qquad
\nu(A\cap[us])>(1-\delta)\,\nu([us]).
\]
Let $T=T_{us,\,u}\colon[us]\to[u]$, the degree-$(+1)$ tail map, with constant
derivative $q$ by Lemma~\ref{lem:rn}. The bad set inside the domain is
\[
[us]\setminus A\quad\text{of measure}\;<\delta\,\nu([us]),
\qquad
T^{-1}\bigl([u]\setminus A\bigr)\quad\text{of measure}\;
=\tfrac1q\,\nu([u]\setminus A)<\delta\,\nu([us]),
\]
so $B:=(A\cap[us])\cap T^{-1}(A\cap[u])$ has $\nu(B)>(1-2\delta)\nu([us])>0$, and
$T|_B$ maps a positive-measure subset of $A$ into $A$ with derivative exactly $q$.
Hence $q\in r_\infty$; applying the same argument to $T^{-1}$ and to compositions
gives $q^{\Z}\subseteq r_\infty$.

\emph{The value $0$.} Since $r_\infty\cap\R_{>0}=q^{\Z}\neq\{1\}$, the relation is not
of type II, and for an ergodic relation this forces $0\in r_\infty$ by Krieger's
dichotomy; alternatively, iterating $T_{us^k,\,u}$ realizes derivatives $q^{k}\to\infty$
on shrinking sets, whose inverses realize derivatives $\to0$.
\end{proof}

\subsection{The factor and its type}

\begin{theorem}[Type of the limit]\label{thm:type}
Let $M_\infty:=W^*(\cR,\nu)$ be the Feldman--Moore von Neumann algebra of the
geometric tail relation. Then $M_\infty$ is the amenable (hyperfinite) factor of type
$\mathrm{III}_{1/q}$, isomorphic to the Powers factor $R_{1/q}$. Moreover, for any
quasi-invariant measure class on $\Omega$ whose Radon--Nikodym cocycle on $\cR$ takes
values in a subgroup $\mu_0^{\Z}\subset\R_{>0}$ with $\mu_0\neq1$ and for which the
density argument of Theorem~\ref{thm:ratio} applies verbatim --- in particular for the
Satake-twisted cocycle classes of [I, \S2] whose weights range in $q^{-\beta\Z}$ ---
the factor is of type $\mathrm{III}_{\lambda}$ with
$\lambda=\min(\mu_0,\mu_0^{-1})$; for the twisted classes,
$\mathrm{III}_{q^{-\beta}}$.
\end{theorem}

\begin{proof}
$\cR$ is a countable Borel equivalence relation with countable classes, generated by
the maps $T_{u,v}$, hence amenable as an increasing union of finite subrelations
(the level-$N$ truncations of [I, Def.~1.6]); $\nu$ is quasi-invariant
(Lemma~\ref{lem:rn}) and ergodic (Lemma~\ref{lem:ergodic}), so $M_\infty$ is an
amenable factor by Feldman--Moore. Krieger's theorem identifies the type of an ergodic
amenable relation from its ratio set: $r_\infty\cap\R_{>0}=\lambda^{\Z}$ with
$0<\lambda<1$ gives type $\mathrm{III}_\lambda$. Theorem~\ref{thm:ratio} gives
$\lambda=1/q$. Uniqueness of the amenable type $\mathrm{III}_\lambda$ factor for
$0<\lambda<1$ (Connes) identifies $M_\infty$ with the Araki--Woods/Powers factor
$R_{1/q}$. The twisted statement repeats the proof with $q$ replaced by $\mu_0$
wherever the cocycle value enters; nothing else in the argument refers to the
numerical value.
\end{proof}

\begin{remark}[A correction to the loose phrasing of {[I, Thm.~2.4(iii)]}]
\label{rem:correction}
Paper I stated a family ``type $\mathrm{III}_{q^{-\beta}}$'' for the boundary system.
For the \emph{untwisted} geometric cocycle this is misleading as a family: the uniform
measure is the unique conformal probability class for every $\beta$ (all branch
weights being equal, the normalized cylinder masses are $\beta$-independent), the KMS
point on the boundary quotient is unique, and the factor type is
$\mathrm{III}_{1/q}$, full stop. A genuine $\beta$-family of types arises only when
the cocycle itself is deformed --- the Satake-twisted classes of
Theorem~\ref{thm:type} --- and the correct statement is the cocycle-range statement
proved above. We record this as a correction of emphasis to [I], not of any equation.
\end{remark}

\subsection{Dissolution, and its exact purchase}

\begin{corollary}[Dissolution of formulation]\label{cor:dissolve}
$M_\infty$ admits no faithful normal semifinite trace. Consequently no inequality of
the form $\tau(x^*x)\ge0$ with $\tau$ a trace is available on $M_\infty$, and the
finite-level obstruction mechanism of [II, Thms.~1.5--1.6] --- non-negativity of
$\Tr(\,\cdot\,^2)$ against $\Tr(C^2)=-n(q-1)<0$ --- has no formulation on the weak
closure of the tower. The obstruction is not overcome; the sentence expressing it
ceases to exist.
\end{corollary}

\begin{proof}
Type III means precisely the absence of a nonzero normal semifinite trace; the rest is
the observation that the finite-level argument is a statement about a trace.
\end{proof}

\begin{proposition}[Inductive persistence]\label{prop:inductive}
Let $(\cE_k)_{k\ge0}$ be genuine $C^*$-correspondences over the level algebras $A_k$,
compatible with the tower transfers, such that $[\cE_k]\in KK^0(A_k,A_k)$ equals the
level-$k$ companion class $C_k$ (equivalently, realizes $\Delta_{p,k}$ as its Pimsner
index map). Then no such system exists: already each single level is impossible by
[II, Thm.~1.6], and a fortiori no inductive limit
$\cE_\infty=\varinjlim\cE_k$ realizes $[\mathsf D_p]_\infty$.
\end{proposition}

\begin{proof}
Level by level, [II, Thm.~1.6]; there is nothing to add, which is the point.
\end{proof}

\begin{remark}[What dissolution buys, and what it destroys]\label{rem:honest}
Corollary~\ref{cor:dissolve} and Proposition~\ref{prop:inductive} bracket the truth
from two sides, and a third fact closes the bracket. The von Neumann completion that
dissolves the trace also dissolves the prize: $K_0$ of a type III factor is trivial,
so $M_\infty$ retains no shadow of $\Jac(Y_k)$ --- the weak closure destroys exactly
the invariant the programme exists to realize. The dissolution theorem therefore does
\emph{not} say ``pass to the limit and construct''; it says that the unique
\emph{numerical} obstruction known dies at the measure-theoretic level, while the
$K$-theoretic content survives only at the $C^*$-level. The sharpened open problem,
which we state as the successor to [II, Conj.~3.6(iii)], is accordingly: does there
exist a \emph{traceless, purely infinite, nuclear} $C^*$-algebra over
$A_\infty=C(\Omega)$ --- a Kirchberg-type limit, where no trace exists to obstruct and
where $K$-theory is nontrivial --- carrying a genuine correspondence whose index map
recovers the tower $(\Z\oplus\Jac(Y_k))_k$? Tracelessness removes the only obstruction
we have proved; it does not supply the construction, and we do not possess one.
\end{remark}

\subsection{The \texorpdfstring{$C^*$}{C*}-level: traceless with
\texorpdfstring{$K$}{K}-theory alive}\label{subsec:kirchberg}

\begin{proposition}[$C^*$-dissolution, no $K$-annihilation]\label{prop:kirchberg}
Let $C^*_r(\cG_\infty)$ denote the reduced groupoid $C^*$-algebra of the rooted tail
groupoid of [I, Def.~1.6] over $\Omega$. It is a unital Kirchberg algebra: separable,
nuclear, simple, purely infinite, and in the UCT class. Being purely infinite, it
admits no nonzero trace, finite or semifinite, densely defined or bounded. Its $K$-theory is
\[
K_0\bigl(C^*_r(\cG_\infty)\bigr)\cong\Z/(q-1),\qquad
K_1\bigl(C^*_r(\cG_\infty)\bigr)=0,
\]
by stabilization to the Cuntz algebra $\mathcal O_q$ $($[I, Prop.~1.7(iii)]$)$.
Consequently, at the $C^*$-level the trace obstruction of [II] has no formulation
\emph{and} $K$-theory is not annihilated: both properties denied to the von Neumann
closure of Corollary~\ref{cor:dissolve} and Remark~\ref{rem:honest} hold
simultaneously here.
\end{proposition}

\begin{remark}[What this algebra does not do, stated against temptation]
\label{rem:blindagain}
It would be pleasant to add that $K_0(C^*_r(\cG_\infty))$ ``carries the tower of
sandpile groups.'' It does not, and asserting it would contradict this series' own
results: the boundary channel is arithmetically blind ([I, Thm.~3.7]), its $K$-theory
seeing only $\Z/(q-1)$ locally and Euler-characteristic data globally, and nothing of
$\Jac(Y_k)$. Proposition~\ref{prop:kirchberg} therefore proves an \emph{existence of
habitat}, not an occupancy: the $C^*$-level is the unique regime where tracelessness
and nontrivial $K$-theory coexist, but the particular algebra naturally supplied by
the tail groupoid is the wrong tenant. The sharpened open problem of
Remark~\ref{rem:honest} stands exactly as posed: construct, functorially from the
Hecke data, a Kirchberg algebra over $C(\Omega)$ whose $K_0$ realizes
$(\Z\oplus\Jac(Y_k))_k$; abstract existence for any prescribed UCT $K$-data is
guaranteed by the Kirchberg--Phillips classification (quoted, data-unverified per the
series policy), and canonicity, not existence, is the entire problem.
\end{remark}

\section{The Iwahori--Hecke Euler characteristic}\label{sec:euler}

The congruence (Iwahori path) tower of [II, \S5] has levels
$Y_k$ = non-backtracking $k$-paths of $Y_0$, with $n_k=n(q+1)q^{k-1}$ vertices for
$k\ge1$, connected by the two degeneracy correspondences (initial and terminal
truncation), and it is non-normal: no deck group acts, and [II, Rem.~5.2] showed its
torsion growth has the exact pseudo-Iwasawa form with $\lambda=-1$, impossible for a
$\Z_p[[T]]$-characteristic ideal. This section derives that law from a single
combinatorial input and locates the $-1$.

\begin{theorem}[Line-digraph recursion, proved]\label{thm:recursion}
Let $G$ be a strongly connected digraph on $n$ vertices in which every vertex has
out-degree and in-degree $q\ge2$ (Eulerian $q$-regular), $m=qn$ arcs, and let $L(G)$
be its line digraph. Then $L(G)$ is again strongly connected and Eulerian
$q$-regular on $m$ vertices, and
\[
\bigl|\mathrm{sandpile}(L(G))\bigr|
=q^{\,m-n-1}\,\bigl|\mathrm{sandpile}(G)\bigr|
=q^{\,n(q-1)-1}\,\bigl|\mathrm{sandpile}(G)\bigr| .
\]
In particular, along the Iwahori tower ($k\ge1$),
$|\mathrm{sandpile}(Y_{k+1})|=q^{\,n_k(q-1)-1}|\mathrm{sandpile}(Y_k)|$; for $q=2$
this is the five-level-verified $2^{\,n_k-1}$ of [II, Rem.~5.2], whose exponent is a
$q=2$ coincidence, corrected here.
\end{theorem}

\begin{proof}
\emph{Step 1 (spectra).} Let $D_{\mathrm{out}},D_{\mathrm{in}}\in M_{n\times m}(\Z)$
be the tail and head incidence matrices, $(D_{\mathrm{out}})_{v,a}=[\,\mathrm{tail}(a)
=v\,]$, $(D_{\mathrm{in}})_{v,a}=[\,\mathrm{head}(a)=v\,]$. Then
$A(G)=D_{\mathrm{out}}D_{\mathrm{in}}^{\,t}$ (arcs $u\to v$) and
$A(L(G))=D_{\mathrm{in}}^{\,t}D_{\mathrm{out}}$ (consecutive arc pairs). For
rectangular $P\in M_{n\times m}$, $Q\in M_{m\times n}$ one has
$\det(zI_m-QP)=z^{\,m-n}\det(zI_n-PQ)$; applied with $P=D_{\mathrm{out}}$,
$Q=D_{\mathrm{in}}^{\,t}$,
\begin{equation}\label{eq:abba}
\det\bigl(zI_m-A(L(G))\bigr)=z^{\,m-n}\det\bigl(zI_n-A(G)\bigr):
\end{equation}
the line digraph has the same nonzero spectrum with multiplicity, plus $m-n$ extra
zeros. $L(G)$ is Eulerian $q$-regular because both degrees of an arc equal the
corresponding degrees of its endpoints in $G$, and strongly connected because a path
between two arcs is obtained by concatenating a vertex path of $G$ from the head of
the first to the tail of the second.

\emph{Step 2 (equal cofactors and the eigenvalue product).} Let
$\Delta=qI-A$ for an Eulerian $q$-regular strongly connected digraph on $N$ vertices.
Row sums of $\Delta$ vanish (out-degree $q$) and column sums vanish (in-degree $q$),
so the all-ones vector spans both the right and the left kernel; by Perron--Frobenius,
$q$ is an algebraically simple eigenvalue of the irreducible non-negative matrix $A$,
so $0$ is a simple eigenvalue of $\Delta$. The two-sided kernel $\C\mathbf1$ forces
all $N$ principal cofactors of $\Delta$ to be equal --- cofactors agree along a row
because the columns of $\Delta$ sum to zero, and along a column because the rows do
--- and the common value is by definition $\kappa:=|\mathrm{sandpile}|$, the
determinant of the reduced Laplacian at any sink; sink-independence is exactly this
equality. Expanding $\det(zI-\Delta)=z\prod_{i=2}^{N}(z-\mu_i)$ over the nonzero
eigenvalues $\mu_i$, the coefficient of $z$ equals $(-1)^{N-1}\prod_{i\ge2}\mu_i$ and
also $(-1)^{N-1}\sum(\text{principal cofactors})=(-1)^{N-1}N\kappa$, whence
\begin{equation}\label{eq:prodform}
\prod_{\mu\in\operatorname{Spec}\Delta\setminus\{0\}}\mu\;=\;N\,\kappa .
\end{equation}

\emph{Step 3 (assembly).} Apply \eqref{eq:prodform} to $G$ and to $L(G)$. By
\eqref{eq:abba}, the nonzero eigenvalues of $\Delta_{L(G)}=qI_m-A(L(G))$ are those of
$\Delta_G$ together with $q-0=q$ taken with multiplicity $m-n$; the eigenvalue $0$
remains simple for $\Delta_{L(G)}$ since $q$ is Perron-simple for $A(L(G))$ as well.
Hence
\[
m\,\kappa(L(G))=q^{\,m-n}\cdot n\,\kappa(G)
\quad\Longrightarrow\quad
\kappa(L(G))=\frac{n}{m}\,q^{\,m-n}\kappa(G)=q^{\,m-n-1}\kappa(G),
\]
using $m=qn$. Besides the five $q=2$ tower levels, the formula was confirmed exactly
in a $q=3$ instance (the complete digraph on four vertices: $\kappa=16$,
$\kappa(L)=34992=3^{7}\cdot16$), against which the uncorrected exponent $n-1$ fails by
a factor $3^{4}$.
\end{proof}

\begin{theorem}[The pseudo-Iwasawa law]\label{thm:pseudo}
Let the base graph have $n$ vertices and $q=p$. Then, unconditionally by
Theorem~\ref{thm:recursion},
for $k\ge2$,
\[
\ord_p\bigl|\mathrm{sandpile}(Y_k)\bigr|
=\underbrace{\frac{n(q+1)}{q}}_{\mu}\,p^{\,k}\;-\;k\;+\;\nu_0,
\qquad
\lambda=-1,
\]
with $\nu_0$ determined by the first level. For $Y_0=K_4$, $q=2$:
$\mu=6$, $\nu_0=-4$, reproducing $\nu_2=6\cdot2^k-k-4$ exactly.
\end{theorem}

\begin{proof}
Summing the recursion of Theorem~\ref{thm:recursion} from level $1$,
\[
\ord_p|\mathrm{sandpile}(Y_k)|
=\ord_p|\mathrm{sandpile}(Y_1)|+\sum_{j=1}^{k-1}\bigl(n_j(q-1)-1\bigr),
\]
the multiplier being a pure power of $p=q$, and
$\sum_{j=1}^{k-1}n_j(q-1)=n(q+1)\bigl(q^{k-1}-1\bigr)=\tfrac{n(q+1)}q\,p^{k}-n(q+1)$.
With $q=2$, $n=4$: $\sum_{j=1}^{k-1}(12\cdot2^{j-1}-1)=6\cdot2^{k}-12-(k-1)$, and
$\ord_2|\mathrm{sandpile}(Y_1)|=7$ gives $6\cdot2^k-k-4$. The general coefficients
read off the same computation.
\end{proof}

\begin{proposition}[$\lambda\ge0$ for commutative Iwasawa towers]\label{prop:lampos}
Let $g\in\Z_p[[T]]$ be nonzero with Weierstrass decomposition
$g=p^{\mu}P(T)U(T)$, $P$ distinguished of degree $\lambda\ge0$, $U$ a unit. Then for
$k\gg0$,
$\ord_p\prod_{\zeta^{p^k}=1,\zeta\neq1}g(\zeta-1)=\mu(p^k-1)+\lambda k+O(1)$
with $\lambda\ge0$. Hence no tower governed by the characteristic ideal of a finitely
generated torsion $\Z_p[[T]]$-module can exhibit $\lambda=-1$: Theorem~\ref{thm:pseudo}
is a rigorous diagnostic that the Iwahori tower is not such a tower.
\end{proposition}

\begin{proof}
$\ord_p g(\zeta_{p^j}-1)=\mu+\lambda\,\ord_p(\zeta_{p^j}-1)+O(p^{-j})
=\mu+\lambda/\varphi(p^j)+o(1)$ for $j\gg0$ by the decomposition, and summing over the
$\varphi(p^j)$ primitive roots for $1\le j\le k$ gives the formula; $\lambda$ is a
degree, hence non-negative.
\end{proof}

\begin{definition}[Degeneracy pair and its Euler term]\label{def:degeneracy}
Let $s_k,t_k\colon A_{k+1}\rightrightarrows A_k$ be the initial- and
terminal-truncation correspondences of the tower, the graph-level avatars of the two
degeneracy maps $X_0(p^{k+1})\rightrightarrows X_0(p^k)$, generating the action of the
Iwahori--Hecke algebra $\cH(\PGL_2(\Q_p),\cI)$ on the tower. The \emph{Eisenstein
kernel} at level $k$ is the one-dimensional space of $\cR$-harmonic constants (the
Perron direction), simultaneously the kernel of the level-$k$ sandpile Laplacian and
the sole obstruction to exactness of the pair $(s_k,t_k)$ on $0$-chains.
\end{definition}

\begin{theorem}[The $-1$ is an Euler characteristic]\label{thm:euler}
The transfer of torsion along one level of the Iwahori tower factors, by
Theorem~\ref{thm:recursion}, as
\[
\ord_p\bigl|\mathrm{sandpile}(Y_{k+1})\bigr|-\ord_p\bigl|\mathrm{sandpile}(Y_k)\bigr|
=\underbrace{n_k(q-1)}_{\text{arc growth: }\ord_p q^{\,m_k-n_k}}
\;-\;\underbrace{\chi_k}_{\;=\;1\;},
\]
where $\chi_k=\dim\ker\Delta^{\mathrm{dir}}_{k}=1$ is the rank of the Eisenstein
kernel removed at each level. Hence in the growth law of Theorem~\ref{thm:pseudo} the
invariants separate cleanly: $\mu$ is the arc-proliferation density
$n(q+1)/q$ of the correspondences, and $\lambda=-\chi=-1$ is the Euler-characteristic
correction of the degeneracy pair --- a homological constant of the non-commutative
base $\cH(\PGL_2(\Q_p),\cI)$, not a module invariant over any commutative
$\Z_p[[T]]$. In particular the corrected form of [II, Conj.~3.6(ii-b)] should assert
the following: the characteristic object of the congruence tower is a module over
$\cH(\PGL_2(\Q_p),\cI)$ whose growth law carries the intrinsic correction
$-\chi\cdot k$, and it reduces to a commutative characteristic ideal exactly on
sub-towers with abelian deck action, where the $\chi$-correction vanishes and the
Newton identities of [II, \S5] hold to the integer.
\end{theorem}

\begin{proof}
The recursion's exponent splits as $(m_k-n_k)-1=\ord_p(q^{\,m_k-n_k})-1$, and both
terms are now located inside the proof of Theorem~\ref{thm:recursion} itself. The
first is the $p$-valuation of the extra Perron factors: each of the $m_k-n_k$ new
zero eigenvalues of $A(Y_{k+1})$ contributes $q-0=q$ to the eigenvalue product
\eqref{eq:prodform}. The subtracted $1$ is the vertex-count normalization
$n_k/m_k=q^{-1}$ in \eqref{eq:prodform}: precisely the one-dimensional Eisenstein
kernel $\ker\Delta^{\mathrm{dir}}$, simple at every level by strong connectivity,
excluded from the eigenvalue product once per level, independently of $k$. Summation
as in Theorem~\ref{thm:pseudo} turns the per-level $-\chi_k=-1$ into the global
$-k$, i.e.\ $\lambda=-\chi=-1$.
\end{proof}

\begin{remark}[Honest scope]\label{rem:scope2}
Theorem~\ref{thm:euler} is now unconditional, but its final sentence remains a
proposal for restating
[II, Conj.~3.6(ii-b)], not a theorem about $\cH(\PGL_2(\Q_p),\cI)$-modules: we have
defined no characteristic ideal over the Iwahori--Hecke algebra, and constructing that
module theory --- with a Weierstrass-type preparation in which $-\chi$ appears
intrinsically --- is the open algebraic problem this section poses. What is
unconditional is Proposition~\ref{prop:lampos}: $\lambda=-1$ certifies, rigorously,
that whatever that theory is, it is not commutative Iwasawa theory.
\end{remark}

\section{Exploratory outlook}\label{sec:outlook}

We record, deriving nothing, where the two threads point. The dissolution theorem
directs the realization problem toward traceless Kirchberg limits over $C(\Omega)$
(Remark~\ref{rem:honest}), where the natural analytic invariants are not traces but
the KMS functionals of the twisted dynamics of [I, \S2], with local $L$-factors as
partition functions. The Euler-characteristic correction of
Section~\ref{sec:euler} lives at the trivial zero of those $L$-factors --- the same
Eisenstein direction removed once per level --- and it is natural to ask, in the
spirit of the exceptional-zero phenomena of Mazur--Tate--Teitelbaum and of
Gross--Stark, whether the per-level derivative data of the tower assembles into a
$p$-adic $L$-functional on the Iwahori--Hecke base whose trivial-zero corrections
reproduce $-\chi k$. We state this as a direction, not a conjecture: no formulation
precise enough to be false is offered here, and the series' methodology is to compute
before naming. The first computable instance has in fact been computed. Write
$g(T)=T^2h(T)$ for the cleared voltage pencil (the double trivial zero of [II,
Rem.~5.4]) and
$\varepsilon_j:=\ord_\ell\prod_{\zeta\ \mathrm{prim.}\ \ell^j}h(\zeta-1)
-\bigl(\mu\,\varphi(\ell^j)+\lambda_h\bigr)$
for the exceptional low-level defects. Then the constant term of the growth law
satisfies, exactly in all four abelian towers of [II, Table~2],
\[
\nu_0=\ord_\ell\kappa(Y_0)-\mu+\sum_{j\ge1}\varepsilon_j ,
\]
with $\varepsilon$ finitely supported: identically zero in three towers, and
$(\varepsilon_1,\varepsilon_2)=(-1,+2)$ in the $(1,2,0)$ tower. The trivial-zero
residue $a_2=\tfrac12g''(0)$ carries the leading datum precisely when $\lambda_h=0$,
where $\ord_\ell a_2=\mu$ as in both $(1,0,0)$ towers, and not otherwise: an
exceptional-zero reading in the Mazur--Tate--Teitelbaum style is available only in
the $\lambda_h=0$ regime, and this bookkeeping identity, verified to the integer, is
the whole of what is presently proved.

\section*{Acknowledgement of provenance and caveats}

Unrefereed preprint. Proved in full: Lemmas~\ref{lem:rn}--\ref{lem:ergodic},
Theorem~\ref{thm:ratio}, Corollary~\ref{cor:dissolve},
Propositions~\ref{prop:inductive} and~\ref{prop:lampos},
Theorem~\ref{thm:recursion} (self-contained; the only classical input is the
algebraic simplicity of the Perron eigenvalue of an irreducible non-negative matrix),
Theorem~\ref{thm:pseudo}, and Theorem~\ref{thm:euler}, all now unconditional.
Quoted classical inputs: every entry in the bibliography is verified against its
primary source or Crossref record --- Connes and Powers via JSTOR, Araki--Woods via
EMS Press, Krieger via Springer, Feldman--Moore Parts I and II via Crossref and the
AMS Digital Library, Levine via Crossref --- with no residual soft spots; the DOIs
are recorded inline at each entry. The Levine reference
\cite{LevineIII} is retained for priority context only: Theorem~\ref{thm:recursion} is
proved self-containedly above, and comparing the exponent $q^{\,n(q-1)-1}$ with the
published statement belongs to the primary-source check this series requires. The concluding restatement of [II, Conj.~3.6(ii-b)] in
Theorem~\ref{thm:euler} remains a proposal.
Remark~\ref{rem:correction} records a correction of emphasis to [I, Thm.~2.4(iii)].
Section~\ref{sec:outlook} makes no claims.

\begingroup\small
\begin{thebibliography}{9}
\bibitem{Powers} R.~T.~Powers, \emph{Representations of uniformly hyperfinite
algebras and their associated von Neumann rings}, Ann.\ of Math.\ (2) \textbf{86}
(1967), no.~1, 138--171. [Primary-source verified: JSTOR, DOI 10.2307/1970364.]
\bibitem{ArakiWoods} H.~Araki and E.~J.~Woods, \emph{A classification of factors},
Publ.\ Res.\ Inst.\ Math.\ Sci.\ \textbf{4} (1968), no.~1, 51--130. [Primary-source
verified: EMS Press / PRIMS article page.]
\bibitem{Krieger} W.~Krieger, \emph{On ergodic flows and the isomorphism of factors},
Math.\ Ann.\ \textbf{223} (1976), 19--70. [Primary-source verified: Springer,
Math.\ Ann.\ \textbf{223}, pp.~19--70, February 1976.]
\bibitem{Connes} A.~Connes, \emph{Classification of injective factors. Cases
$\mathrm{II}_1$, $\mathrm{II}_\infty$, $\mathrm{III}_\lambda$, $\lambda\neq1$},
Ann.\ of Math.\ (2) \textbf{104} (1976), no.~1, 73--115. [Primary-source
verified: JSTOR, DOI 10.2307/1971057.]
\bibitem{FeldmanMoore} J.~Feldman and C.~C.~Moore, \emph{Ergodic equivalence
relations, cohomology, and von Neumann algebras.\ I, II}, Trans.\ Amer.\ Math.\ Soc.\
\textbf{234} (1977), 289--324 and 325--359. [Both parts verified: Part I
Crossref-verified, DOI 10.1090/S0002-9947-1977-0578656-4, pp.~289--324; Part II
primary-source verified, DOI 10.1090/S0002-9947-1977-0578730-2, pp.~325--359. The
related 1975 Bulletin announcement, Bull.\ Amer.\ Math.\ Soc.\ \textbf{81} (1975),
no.~5, 921--924, DOI 10.1090/S0002-9904-1975-13888-2, is a distinct paper and is not
the reference intended here.]
\bibitem{LevineIII} L.~Levine, \emph{Sandpile groups and spanning trees of directed
line graphs}, J.\ Combin.\ Theory Ser.\ A \textbf{118} (2011), no.~2, 350--364.
[Crossref-verified: DOI 10.1016/j.jcta.2010.04.001.]
\end{thebibliography}
\endgroup

\end{document}
